% Coverage: physical pp. 98--100, Section 16.4 from its heading through Eq. (16.4.12), stopping before Problems; all numbered and unnumbered displays included.
% Uncertainties: none after rendered-page review; Grassmann derivative sides and signs were checked directly against physical p. 100.
\subsection{Symmetries of the Effective Action}

In some but not all cases, the symmetries of the action $I[\phi]$ are automatically also symmetries of the effective action $\Gamma[\phi]$. For instance, in the example of Section 16.2, the action \eqref{eq:16.2.1} has a symmetry under the discrete transformation $\phi\to-\phi$. Thus it follows from their definition that $Z[J]$ and $W[J]$ are even under the corresponding reflection $J\to-J$. Eq.~\eqref{eq:16.1.5} then shows that $\phi_{-J}=-\phi_J$, and hence $J_{-\phi}=-J_\phi$, so Eq.~\eqref{eq:16.1.6} shows that $\Gamma[\phi]$ is even under $\phi\to-\phi$. This is borne out by the one-loop result \eqref{eq:16.2.15}. The fermion-loop contribution in \eqref{eq:16.2.16} also exhibits the symmetry under $\phi\to-\phi$ in the special case where $M(0)=0$, because in this case the action is invariant under the combined transformation $\phi\to-\phi$, $\psi\to\gamma^5\psi$.

We encounter problems in establishing the renormalizability of a theory, unless we can show that the symmetries that we impose on the action also apply to the effective action. For instance, in the example above, if $I[\phi]$ were assumed to be even in $\phi$ but $\Gamma[\phi]$ turned out not to be, then the coefficients of the terms in $\Gamma$ proportional to $\int d^4x\,\phi$ and $\int d^4x\,\phi^3$ would be divergent, but the symmetry of the action would not allow us to introduce counterterms to absorb these infinities.

With this motivation, let's now turn to the important class of symmetries generated by infinitesimal transformations
\begin{equation}
\chi^n(x)\longrightarrow\chi^n(x)+\epsilon F^n[x;\chi],
\tag{16.4.1}\label{eq:16.4.1}
\end{equation}
where $F^n$ is a function of $x^\mu$ that depends functionally on $\chi^n$. (For instance, $F^n[x;\chi]$ may be an ordinary function of the $\chi^n$ and their derivatives at the point $x$.) We are now using the symbol $\chi^n$ rather than $\phi^n$ to denote the different types of fields, to emphasize that these include not only ordinary gauge and matter fields (which in the next chapter will be denoted $\phi^r(x)$), but all other fields appearing in the gauge-fixed action, including ghost fields. We repeat that these $\chi^n(x)$ may be of any type, not necessarily scalars.

We assume that both the action and the measure are invariant under the symmetry transformation \eqref{eq:16.4.1}:
\begin{align}
I[\chi+\epsilon F]&=I[\chi],
&&\tag{16.4.2}\label{eq:16.4.2}\\
\prod_{n,x}d\bigl(\chi^n(x)+\epsilon F^n[x;\chi]\bigr)
&=\prod_{n,x}d\chi^n(x).
&&\tag{16.4.3}\label{eq:16.4.3}
\end{align}
(It is actually sufficient for only the product $(\prod_{n,x}d\chi^n(x))\exp(iI)$ to be invariant, but where this is true usually Eqs.~\eqref{eq:16.4.2} and \eqref{eq:16.4.3} both apply.) Replacing the integration variables in Eq.~\eqref{eq:16.1.1} with $\chi^n(x)+\epsilon F^n[x;\chi]$, we have then
\begin{align*}
Z[J]={}&\int\left[\prod_{n,x}d\bigl(\chi^n(x)+\epsilon F^n[x;\chi]\bigr)\right]
\\[-2pt]
&\quad\cdot\exp\left\{iI[\chi+\epsilon F]
+i\int d^4x\,\bigl(\chi^n(x)+\epsilon F^n[x;\chi]\bigr)J_n(x)\right\}
\\
={}&\int\left[\prod_{n,x}d\chi^n(x)\right]
\exp\left\{iI[\chi]
+i\int d^4x\,\bigl(\chi^n(x)+\epsilon F^n[x;\chi]\bigr)J_n(x)\right\}
\\
={}&Z[J]+i\epsilon\int\left(\prod_{n,x}d\chi^n(x)\right)
\int F^n[y;\chi]J_n(y)\,d^4y
\\[-2pt]
&\quad\cdot\exp\left\{iI[\chi]+i\int d^4x\,\chi^n(x)J_n(x)\right\}
\end{align*}
and hence
\begin{equation}
\int d^4y\,\langle F^n(y)\rangle_J J_n(y)=0,
\tag{16.4.4}\label{eq:16.4.4}
\end{equation}
where $\langle\ \rangle_J$ denotes the quantum average in the presence of the current $J_n(x)$,
\begin{align}
Z[J]\langle F^n(y)\rangle_J\equiv{}&
\int\left(\prod_{n,x}d\chi^n(x)\right)F^n[y;\chi]
\notag\\[-2pt]
&\quad\cdot\exp\left\{iI[\chi]+i\int d^4x\,\chi^n(x)J_n(x)\right\},
\tag{16.4.5}\label{eq:16.4.5}
\end{align}
normalized so that $\langle1\rangle_J=1$. But recall that $J_n(y)$ is given in terms of the effective action $\Gamma[\chi]$ by Eq.~\eqref{eq:16.1.7}
\[
J_{n,\chi}(y)=-\frac{\delta\Gamma[\chi]}{\delta\chi^n(y)}.
\]
Therefore Eq.~\eqref{eq:16.4.4} may be written as
\begin{equation}
0=\int d^4y\,\langle F^n(y)\rangle_{J_\chi}
\frac{\delta\Gamma[\chi]}{\delta\chi^n(y)}.
\tag{16.4.6}\label{eq:16.4.6}
\end{equation}
In other words, $\Gamma[\chi]$ is invariant under the infinitesimal transformation
\begin{equation}
\chi^n(y)\longrightarrow\chi^n(y)+\epsilon\langle F^n(y)\rangle_{J_\chi}.
\tag{16.4.7}\label{eq:16.4.7}
\end{equation}
Such symmetry conditions are known as Slavnov--Taylor identities~\hyperref[ch16-ref-7]{[7]}.

Is this the symmetry transformation with which we started? It is for one very important class of infinitesimal symmetry transformations: those that are \emph{linear}. For such symmetries $F$ is
\begin{equation}
F^n[x;\chi]=s^n(x)+\int t^n{}_m(x,y)\chi^m(y)\,d^4y.
\tag{16.4.8}\label{eq:16.4.8}
\end{equation}
(In the most common case $s^n(x)$ vanishes and $t^n{}_m(x,y)$ is a constant matrix times $\delta^4(x-y)$.) For any linear $F$, we have
\[
\langle F^n(x)\rangle_J
=s^n(x)+\int t^n{}_m(x,y)\langle\chi^m(y)\rangle_J\,d^4y.
\]
But for any fixed $\chi$, $J_\chi$ is defined as the value of the current $J$ that makes $\langle\chi^m(y)\rangle_J$ equal to $\chi^m(y)$, so
\begin{equation}
\langle F^n(x)\rangle_{J_\chi}
=s^n(x)+\int t^n{}_m(x,y)\chi^m(y)\,d^4y
=F^n[x;\chi].
\tag{16.4.9}\label{eq:16.4.9}
\end{equation}
Hence Eq.~\eqref{eq:16.4.6} requires that $\Gamma[\chi]$ be invariant under all the functional linear transformations $\chi^n\to\chi^n+\epsilon F^n$ that leave $I[\chi]$ and the measure invariant.

We occasionally have to deal with symmetry transformations that are not linear. One important example is provided by the BRST transformation discussed in Section 15.7. For non-linear transformations, the symmetry transformation \eqref{eq:16.4.7} under which the effective action is invariant is not generally the same as the assumed symmetry transformation \eqref{eq:16.4.1} that leaves the original action invariant, because the average of a non-linear functional of fields is not generally the same as the functional of the average fields. Indeed, the form of $\langle F\rangle_{J_\chi}$ as a functional of $\chi$ depends in general on the dynamics of the system, and is usually non-local. This complication will be dealt with in the next chapter by the method of antibrackets.

\begin{center}
$*\ *\ *$
\end{center}

Up to now, we have tacitly been assuming that the fields $\chi^n$ and the corresponding transformation functions $\Delta^n$ and currents $J_n$ are all bosonic. We will need to take note of the sign factors that appear when some of these are fermionic, as in particular for supersymmetry or BRST transformations, where $\epsilon$ is fermionic and $\chi^n$ and $\Delta^n$ have opposite statistics. With currents inserted to the right of fields, as in Eqs.~\eqref{eq:16.1.1} and \eqref{eq:16.4.5}, Eqs.~\eqref{eq:16.1.5} and \eqref{eq:16.1.7} hold in the form
\begin{align}
\frac{\delta_R W[J]}{\delta J_m(y)}&=\chi_J^m(y),
&&\tag{16.4.10}\label{eq:16.4.10}\\
\frac{\delta_L\Gamma[\chi]}{\delta\chi^m(y)}&=-J_{\chi,m}(y),
&&\tag{16.4.11}\label{eq:16.4.11}
\end{align}
with the subscripts $R$ and $L$ indicating that the derivative is to act from the right or left. In consequence, the Slavnov--Taylor identity \eqref{eq:16.4.6} should be written
\begin{equation}
0=\int d^4y\,\langle F^n(y)\rangle_{J_\chi}
\frac{\delta_L\Gamma[\chi]}{\delta\chi^n(y)}.
\tag{16.4.12}\label{eq:16.4.12}
\end{equation}
